\chapter{Logica Proposizionale}

Abbiamo visto, nel capitolo precedente, che una logica è caratterizzata da tre componenti fondamentali: un linguaggio formale, una semantica e un calcolo.
In questo capitolo, studieremo in dettaglio la \textit{logica proposizionale}, la più semplice tra le logiche classiche, che costituisce il punto di partenza per comprendere sistemi logici più complessi.
Procederemo sistematicamente, partendo dalla definizione del linguaggio (sintassi), per poi introdurre la semantica e infine studiare i sistemi di calcolo per la derivazione di conseguenze logiche.

\section{Sintassi della Logica Proposizionale}

Il primo passo nella costruzione di una logica formale è definire il suo linguaggio, anche detto alfabeto. Il linguaggio è l'insieme dei simboli da cui attingere per la creazione di espressioni.
In generale, un linguaggio logico è composto da due tipi di simboli: i \textit{simboli non logici}, su cui la logica non fa alcuna assunzione e il cui significato varia a seconda della \qi{situazione}, e i \textit{simboli logici}, i cui significati sono fissati dalla specifica logica.

\begin{definition}{Alfabeto della logica proposizionale}
  L'alfabeto della logica proposizionale è costituito da:
  \begin{itemize}
    \item Un insieme di \textit{simboli non logici}, detti \keyword{proposizioni atomiche} e denotati con \(\AP = \set{A,B,C,\ldots}\).
    \item Un insieme di \textit{simboli logici}, tra cui i connettivi logici \(\set{\land, \lor, \rightarrow, \neg,\leftrightarrow}\) e le parentesi \(\set{(,)}\) per disambiguare la struttura delle formule.
  \end{itemize}
\end{definition}

\begin{note}{}
L'insieme dei simboli non logici è potenzialmente infinito.
\end{note}

Un'\keyword{espressione} è una sequenza finita di simboli del linguaggio.
Tuttavia, non tutte le espressioni sono significative dal punto di vista logico: vogliamo considerare solo quelle che rispettano le regole grammaticali del linguaggio.

\begin{definition}{Formula ben formata (\(\WFF\))}
  Una formula ben formata (\(\WFF\)) è un'espressione che può essere costruita usando una grammatica definita sui simboli del linguaggio.
\end{definition}

Nel caso specifico della logica proposizionale, l'insieme \(\WFF\) viene indicato con \(PL\), che sta per \textit{Propositional Logic}.
In generale, l'insieme delle formule ben formate di una logica viene indicato con le iniziali di tale logica. 

La definizione precisa di quali espressioni siano formule ben formate nella logica proposizionale è data ricorsivamente attraverso le seguenti regole di formazione.

\begin{definition}{\(\WFF\) Proposizionale}
  Nella logica proposizionale, una formula è ben formata se è costruita secondo le seguenti regole:
\begin{itemize}
  \item \(\AP \subseteq \WFF\)
  \item Se \(\phi \in \WFF\), allora \(\neg \phi \in \WFF\)
  \item Se \(\phi,\psi \in \WFF\), allora \((\phi \odot \psi)\in \WFF\), dove \(\odot \in \set{\land, \lor, \rightarrow, \leftrightarrow}\) 
  \item Nessuna altra espressione è una formula ben formata
\end{itemize}
  
\end{definition}

Come si può notare, la definizione è una grammatica libera da contesto, e, di conseguenza, ogni formula ben formata può essere rappresentata come un albero sintattico, in cui i nodi interni corrispondono a connettivi logici e le foglie a proposizioni atomiche.
Tale albero sintattico cattura la struttura gerarchica della formula, mostrando come essa sia stata costruita a partire dalle sue componenti più semplici.

\begin{example}{Albero sintattico}
  Consideriamo la formula \(\phi = (A \rightarrow (B \land \neg C))\). 
  
  Secondo la definizione ricorsiva, questa formula è ben formata: a partire dalle variabili proposizionali \(A\), \(B\), \(C\) (che, per la prima regola, sono \(\WFF\)), costruiamo \(\neg C\) (seconda regola), poi \((B \land \neg C)\) (terza regola), e infine \((A \rightarrow (B \land \neg C))\) (terza regola).
  
  L'albero sintattico di \(\phi\) rappresenta questa costruzione gerarchica:
  
  \begin{center}
    \begin{forest}
      [\(\rightarrow\)
        [\(A\)]
        [\(\land\)
          [\(B\)]
          [\(\neg\)
            [\(C\)]
          ]
        ]
      ]
    \end{forest}
  \end{center}
  
  In questo albero:
  \begin{itemize}
    \item Le foglie sono le variabili proposizionali \(A\), \(B\), \(C\)
    \item I nodi interni sono i connettivi logici \(\rightarrow\), \(\land\), \(\neg\)
    \item La struttura dell'albero riflette l'ordine di applicazione delle regole di costruzione delle \(\WFF\)
  \end{itemize}
\end{example}

Tale costruzione ci permette di utilizzare un potente strumento di ragionamento: l'\textit{induzione strutturale}.
L'induzione strutturale consente di dimostrare proprietà delle formule procedendo per casi: il caso base verifica che la proprietà in esame sia valida per le formule atomiche (variabili proposizionali), mentre i passi induttivi verificano la proprietà su formule costruite a partire da sottoformule più semplici.

A differenza dell'induzione matematica semplice sui numeri naturali, dove per dimostrare \(P(n)\) assumiamo solo \(P(n-1)\), l'induzione strutturale è basata sul principio di \textit{induzione completa} (o forte), secondo cui, per dimostrare una proprietà \(P(n)\) assumiamo che valga \(P(m), \forall m < n\). 
Nel caso dell'induzione strutturale, per dimostrare una proprietà di \(\phi\), assumiamo che tale proprietà valga per ogni sottoformula propria di \(\phi\).
L'induzione strutturale su una formula è equivalente all'induzione completa sull'altezza del relativo albero sintattico. Poiché l'altezza di un albero sintattico fornisce un ordine ben fondato\footnote{L'altezza di un albero indica il numero massimo di nodi lungo un percorso dalla radice a una foglia. L'altezza è un intero non negativo: il minimo è \(0\), ovvero l'altezza di una foglia. Essendo l'insieme degli interi non negativi, assieme alla relazione d'ordine \(\leq\), ben fondato, anche l'insieme delle altezze, con la stessa relazione d'ordine, è ben fondato.}, anche l'induzione strutturale è ben fondata.

La definizione ricorsiva di formula ben formata impone che i sottoalberi di un albero sintattico siano a loro volta alberi sintattici.
Tali sottoalberi rappresentano le \keyword{sottoformule} di una formula, il cui insieme è definito come segue.

\begin{definition}{Sottoformule}
  Dato \(\phi \in WFF\), l'\keyword{insieme delle sottoformule} di \(\phi\), denotato con \(\Sub{\phi}\), è definito ricorsivamente come segue:
  \[\Sub{\phi} = \begin{cases}
    \{\phi\} & ,\phi \in \AP\\
    \Sub{\phi'} \cup \{\neg \phi'\} & ,\phi = \neg \phi'\\
    \Sub{\phi'} \cup \Sub{\psi'} \cup \{(\phi' \odot \psi')\} & ,\phi = (\phi' \odot \psi')
  \end{cases}
  \]
\end{definition}

Possiamo ora osservare una prima applicazione dell'induzione strutturale, dimostrando una proprietà fondamentale delle formule ben formate: il numero di sottoformule di una formula \(\phi\) è al più pari alla lunghezza della formula stessa.

\begin{note}{}
  Tale proprietà evidenzia l'importanza delle regole di formazione delle WFF, differenziandole da generiche successioni di simboli.
  Infatti, data una stringa qualsiasi, il numero di sottostringhe è quadratico rispetto alla lunghezza della stringa\footnote{Si prenda ad esempio la formula \(ABC\), che non è ben formata. In essa, troviamo una sola sottostringa di lunghezza \(3\) (\(ABC\)), due sottostringhe di lunghezza \(2\) (\(AB,BC\)), tre sottostringhe di lunghezza \(1\) (\(A,B,C\)). Si vede quindi che, per una stringa qualsiasi di lunghezza \(n\), il numero di sue sottoformule è \(n(n-1)/2=\mathcal{O}(n^{2})\).}, mentre, presa una WFF, il numero di sottoformule è solamente lineare rispetto alla lunghezza della formula.
\end{note}

\begin{proposition}[namedlabel={prop:subformula-upperbound}{Limite superiore al numero di sottoformule}]{Limite superiore al numero di sottoformule}
  Sia \(\phi \in WFF\). Siano inoltre \(\abs{\phi}\) la lunghezza della formula \(\phi\) in numero di simboli e \(\abs{\Sub{\phi}}\) la cardinalità dell'insieme delle sottoformule di \(\phi\). Allora
  \[\abs{\Sub{\phi}}\leq \abs{\phi}\]
\end{proposition}

\begin{proof}
  Per dimostrare questa proposizione, procediamo per induzione strutturale sulla formula \(\phi\).
  \begin{description}
    \item[\(\phi \in \AP\)]: In questo caso, \(\phi\) è una formula atomica, quindi \(\Sub{\phi} = \{\phi\}\) e \(\abs{\Sub{\phi}} = 1\). Dato che \(\abs{\phi} = 1\) (ogni simbolo ha lunghezza 1), abbiamo
          \[\abs{\Sub{\phi}} = 1 \leq 1 = \abs{\phi}\]
    \item[\(\phi = \neg \phi'\)]: In questo caso, per induzione strutturale, assumiamo che \(\abs{\Sub{\phi'}} \leq \abs{\phi'}\). Dato che \(\Sub{\phi} = \Sub{\phi'} \cup \{\neg \phi'\}\), abbiamo
          \[\abs{\Sub{\phi}} = \abs{\Sub{\phi'}} + 1\]
          Inoltre, \(\abs{\phi} = \abs{\phi'} + 1\) (perché aggiungiamo il simbolo di negazione). Quindi,
          \[\abs{\Sub{\phi}} = \abs{\Sub{\phi'}} + 1 \leq \abs{\phi'} + 1 = \abs{\phi}\]
    \item[\(\phi = (\phi' \odot \psi')\)]: Infine, sempre per induzione strutturale, assumiamo che \(\abs{\Sub{\phi'}} \leq \abs{\phi'}\) e \(\abs{\Sub{\psi'}} \leq \abs{\psi'}\). Dato che \(\Sub{\phi} = \Sub{\phi'} \cup \Sub{\psi'} \cup \{(\phi' \odot \psi')\}\), abbiamo
          \[\abs{\Sub{\phi}} \leq \abs{\Sub{\phi'}} + \abs{\Sub{\psi'}} + 1\]
          In questo caso non è detto valga l'uguaglianza, perché è possibile che \(\Sub{\phi'}\) e \(\Sub{\psi'}\) non siano disgiunti.
          Inoltre, \(\abs{\phi} = \abs{\phi'} + \abs{\psi'} + 1\) (perché aggiungiamo il connettivo). Quindi partendo da:
          \[\abs{\Sub{\phi'}} \leq \abs{\phi'}\]
          sommiamo \(\abs{\Sub{\psi'}}\) ad ambo i membri dell'disequazione:
          \[\abs{\Sub{\phi'}} +\abs{\Sub{\psi'}} \leq \abs{\phi'} + \abs{\Sub{\psi'}} \leq \abs{\phi'} + \abs{\psi'}\]
          e dunque
          \[\abs{\Sub{\phi}} \leq \abs{\Sub{\phi'}} + \abs{\Sub{\psi'}} + 1 \leq \abs{\phi'} + \abs{\psi'} + 1 = \abs{\phi}\]
  \end{description}
\end{proof}

\section{Semantica Proposizionale tramite Modelli}

Definita la sintassi della logica proposizionale, diamo una definizione formale di semantica. La \keyword{semantica} della logica proposizionale formalizza il concetto intuitivo di \q{situazione}, che determina quali proposizioni elementari sono vere e quali false.
Un modo per definire la semantica è attraverso il concetto di \keyword{modello}.
Un modello, in generale, è una struttura matematica che interpreta i simboli del linguaggio e stabilisce le condizioni di verità per le formule.
Nel caso della logica proposizionale, una possibile implementazione di modello è mediante insiemi, per cui un modello è un qualsiasi sottoinsieme di \(\AP\).

\begin{definition}{Modello per la logica proposizionale}
  In logica proposizionale, un modello \(m\) è un qualsiasi sottoinsieme di \(\AP\). In altre parole, \(m \in 2^{\AP}\).\footnote{\(2^{\AP}\) non è altro che una notazione alternativa per l'insieme di tutti i sottoinsiemi di \(\AP\), chiamato anche insieme potenza indicato anche con \(\mathcal{P}(\AP)\).}
\end{definition}

Essendo la logica proposizionale molto semplice, per darle significato è sufficiente una struttura altrettanto semplice: un insieme di proposizioni atomiche che consideriamo vere.
Infatti, un modello non è altro che un insieme di proposizioni elementari tale per cui una proposizione è vera se e soltanto se appartiene al modello.
Affinché un modello sia ben definito su formule più complesse, è necessaria una regola ricorsiva per il soddisfacimento.

\begin{note}{}
Essendo \(\AP\) potenzialmente infinito, anche i modelli possono essere infiniti. Da un punto di vista matematico non è particolarmente importante, ma modelli infiniti portano a problemi computazionali di rappresentazione e manipolazione.
\end{note}

\begin{definition}{Soddisfacimento di una formula da parte di un modello}
  Dato un modello \(m \in 2^{\AP}\) e una formula \(\phi \in \PL\), definiamo la \keyword{relazione di soddisfazione}, \(\models \subseteq 2^{\AP} \times \PL\), come segue:

  \(m\models \phi \) se e solo se vale uno dei seguenti casi:
  \begin{enumerate}
    \item \(\phi \in m\) se \(\phi \in \AP\).
    \item \(m \not\models \phi'\) se \(\phi = \neg \phi'\).
    \item \(m \models \phi_1\) e \(m \models \phi_2\) se \(\phi = \phi_1 \land \phi_2\).
    \item \(m \models \phi_1\) o \(m \models \phi_2\) se \(\phi = \phi_1 \lor \phi_2\).
    \item Se \(m \models \phi_1\) allora \(m \models \phi_2\) se \(\phi = \phi_1 \rightarrow \phi_2\).
    \item \(m \models \phi_1\) se e solo se \(m \models \phi_2\) se \(\phi = \phi_1 \leftrightarrow \phi_2\).
  \end{enumerate}
\end{definition}

Tale definizione rappresenta una procedura ricorsiva per determinare se una formula è soddisfatta da un modello, partendo dalle proposizioni atomiche e risalendo l'albero sintattico della formula.

Il problema decisionale di verificare se \(m \models \phi\) è detto \keyword{model checking}, e, a patto di risolvere il problema di rappresentazione di modelli potenzialmente infiniti, è facilmente decidibile\footnote{Un problema è decidibile quando per esso esiste un algoritmo, la cui terminazione è garantita, che lo risolve. In particolare, un problema è facilmente decidibile quando esiste un tale algoritmo che abbia complessità polinomiale.}, poiché equivale a un'esplorazione in post-ordine dell'albero sintattico di \(\phi\), che ha un numero di nodi pari al numero di sottoformule di \(\phi\), che è al più pari alla lunghezza di \(\phi\).

Avendo finalmente definito la semantica possiamo andare a formalizzare le definizioni informali date nel primo capitolo.
\begin{definition}{Soddisfacibilità}
  Una formula \(\phi\) è \keyword{soddisfacibile} (\(\Sat[\phi]\) o \(\phi \in \Sat\)) se e soltanto se esiste un modello \(m\in 2^{\AP}\) tale che \(m \models \phi\).
\end{definition}

Il problema decisionale di verificare \(\Sat[\phi]\) è detto \keyword{soddisfacibilità} ed è un problema centrale in logica e informatica teorica, poiché molti problemi di interesse, come pianificazione delle azioni, colorazione dei grafi, etc, possono essere ridotti a esso. 
Ciò significa che, se riuscissimo a risolvere il problema di soddisfacibilità in modo efficiente, potremmo risolvere anche molti altri problemi in modo efficiente.
Rispetto al model checking, il problema di soddisfacibilità è spesso più complesso, poiché richiede l'esplorazione dello spazio dei modelli, che cresce esponenzialmente con il numero di proposizioni atomiche coinvolte nella formula\footnote{Se una formula è composta da una sola proposizione elementare \(A\), allora sono necessari due modelli per esplorare tutte le situazioni possibili: uno in cui \(A\) è vera, uno in cui \(A\) è falsa. Se \(\phi\) è composta da due formule elementari \(A\) e \(B\), allora, per esplorare tutte le situazioni possibili, ovvero tutte le combinazioni di verità e falsità delle di \(A\) e \(B\), sono necessari quattro modelli. Se le formule sono tre, servono otto modelli e così via. In generale, per descrivere tutti casi possibili per una formula con \(n\) proposizioni elementari, sono necessari \(2^{n}\) modelli.}. 
Più semplicemente, il problema di soddisfacibilità è più complesso perché costituito da tanti problemi di model checking per quanti sono i modelli possibili.

\begin{definition}{Validità}
  Una formula \(\phi\) è \keyword{valida} (\(\Valid[\phi]\) o \(\phi \in \Valid\)) se e soltanto se per ogni modello \(m\in 2^{\AP}\) vale \(m \models \phi\).
\end{definition}

La validità e la soddisfacibilità sono concetti duali, nel senso che \(\Sat[\phi]\iff \neg \Valid[\neg \phi]\) e \(\Valid[\phi] \iff \neg \Sat[\neg \phi]\).\footnote{I simboli \(\iff\) e \(\implies\) sono simboli di \textbf{metalinguaggio}. Tali simboli non appartengono al linguaggio oggetto del discorso (in questo caso la logica proposizionale), ma sono simboli sono utilizzati (con l'interpretazione usuale) per parlare di proprietà del linguaggio stesso.}
Questo è diretta conseguenza del fatto che, data una qualsiasi proprietà \(P\), \(\exists x:P(x) \iff \forall x, \neg P(\neg x)\). Ciò verrà dimostrato formalmente quando tratteremo la logica del primo ordine.

\begin{definition}{Estensione della soddisfacibilità}
  Dato \(\Gamma \subseteq \PL\) e \(m \in 2^{\AP}\), diremo che \(m\) soddisfa \(\Gamma\) (\(m \models \Gamma\)) se e soltanto se \(m \models \phi\) per ogni \(\phi \in \Gamma\).
\end{definition}
Ciò consente una scrittura più compatta per intendere che un modello soddisfa un insieme di formule.

\begin{definition}{Conseguenza logica}
  Dato \(\Gamma \subseteq \PL\), una formula \(\phi\) è conseguenza logica di \(\Gamma\) (\(\Gamma \models \phi\)) se e soltanto se per ogni \(m \in 2^{\AP}\) che soddisfa \(\Gamma\) allora \(m \models \phi\).

  Alternativamente, \(\Gamma \models \phi\) se e soltanto se non esiste \(m \in 2^{\AP}\) tale che \(m \models \Gamma\) e \(m \not\models \phi\).
\end{definition}

\begin{note}{}
  Nonostante si usi la stessa notazione \(\models\) per indicare sia la soddisfacibilità che la conseguenza logica, è importante sottolineare che si tratta di due concetti distinti, anche se tra loro molto legati.
  La soddisfacibilità \(m \models \phi\) è una relazione contenuta in \(2^{\AP}\times \PL\), mentre la conseguenza logica \(\Gamma \models \phi\) è una relazione contenuta in \(2^{\PL}\times \PL\).
\end{note}

Dalla definizione di conseguenza logica segue che se \(\emptyset \models \phi\), denotato alternativamente con \(\models \phi\), allora \(\phi\in\Valid\).
Se consideriamo \(\Gamma\) come l'insieme dei vincoli a cui il modello è sottoposto, allora è chiaro che, se \(\Gamma\) è vuoto, e quindi non vi sono vincoli, allora ogni modello lo soddisfa.
Quindi, la relazione di conseguenza logica dall'insieme vuoto si riduce alla soddisfacibilità di \(\phi\) da parte di ogni modello, ovvero la validità di \(\phi\).
In altre parole, potendo risolvere il problema decisionale di conseguenza logica, potremmo risolvere quello di validità.

Inoltre, se \(\Gamma\) è un insieme finito di formule, si verifica anche il viceversa, ovvero è possibile risolvere il problema decisionale di validità per risolvere quello di conseguenza logica. Infatti, in tal caso, è possibile esprimere \(\Gamma \models \phi\), ovvero il concetto di conseguenza logica di \(\phi\) rispetto a \(\Gamma\), in termini di validità della formula
\[\left(\bigwedge_{\gamma \in \Gamma} \gamma\right) \rightarrow \phi.\]

Ciò, oltre a evidenziare la stretta connessione tra conseguenza logica e validità (e quindi soddisfacibilità, per dualità), sottolinea come la relazione \(\models\) sia correlata alla semantica del simbolo di implicazione \(\rightarrow\), cosa facilmente evincibile dall'analogia tra definizione di conseguenza logica e di semantica dell'implicazione.

\section{Semantica Proposizionale tramite Assegnamenti}

Definito formalmente un primo concetto di semantica, andiamo a risolvere il problema della sua rappresentabilità definendo un altro tipo di semantica, che dimostreremo essere equivalente a quella precedente.

La semantica mediante modelli è semplice ed elegante. 
Tuttavia, data una formula soddisfacibile, esistono infiniti modelli che la soddisfano: un modello potrebbe contenere anche proposizioni elementari non contenute in tale formula.
L'infinitezza dei modelli è un problema non trascurabile per la computazione.
Per risolverlo, possiamo semplicemente limitare la semantica alle sole proposizioni atomiche che compaiono nella formula d'interesse.

\begin{definition}{Proposizioni atomiche di una formula}
  Data \(\phi \in \PL\), possiamo definire \(\AP[\phi] = \Sub{\phi} \cap \AP\) l'insieme delle proposizioni atomiche contenute in \(\phi\).
\end{definition}

Essendo \(\abs{\Sub{\phi}}\leq \abs{\phi}\), si ha \(\abs{\AP[\phi]}\leq \abs{\phi}\). Essendo \(\abs{\phi}\) finito, anche \(\abs{\AP[\phi]}\) lo sarà.

\begin{definition}{Assegnamento}
  Dato \(\phi \in \PL\), un assegnamento \(\alpha\) per \(\phi\) è una funzione \(\alpha: \Delta_\phi \rightarrow \{0,1\}\), tale che \(\Delta_\phi\) è finito e \(\AP[\phi] \subseteq \Delta_\phi \subseteq \AP\).\footnote{Il dominio \(\Delta_{\phi}\) di un assegnamento è un qualsiasi insieme di proposizioni che contenga le proposizioni elementari di \(\phi\). Un assegnamento per cui \(\Delta_{\phi}=\AP[\phi]\) è un assegnamento minimale, ovvero un assegnamento che copre solo le proposizioni strettamente necessarie.}

  Chiameremo l'insieme di tali assegnamenti \(\Ass{\phi}\).\footnote{Analogamente alla notazione \(2^S\) per indicare l'insieme \(\mathcal{P}(S)\), una notazione molto in uso per indicare un insieme di funzioni \(f: D \to C\) è \(C^D\) essendo che la cardinalità di tale insieme (per insiemi finiti) è proprio \(\abs{C}^{\abs{D}}\). Di conseguenza, una notazione alternativa per \(\Ass{\phi}\) è \(\set{0,1}^{\Delta_\phi}\)}
\end{definition}

Come visto nella definizione, data una formula \(\phi\), l'insieme degli assegnamenti compatibili è finito.
In particolare, essi ne sono \(\abs{\set{0,1}}^{\abs{\AP[\phi]}}=2^{\abs{\AP[\phi]}}\).
Tali funzioni di assegnamento ci permettono di assegnare il valore di verità a ciascuna proposizione atomica considerata. 
Tuttavia, per calcolare il valore di verità di una formula complessa, definiamo ora la sua \keyword{valutazione} in modo ricorsivo, a partire dall'assegnamento, un po' come avevamo fatto per la soddisfacibilità nella semantica tramite modelli.

\begin{definition}{Valutazione di una formula}
  Siano \(\phi \in \PL\) e \(\alpha \in \Ass{\phi}\). Definiamo dunque la \keyword{valutazione basata su \(\alpha\)} come una funzione \(\Val{\alpha}: \PL \to \set{0,1}\) come\footnote{Per essere precisi, \(\Val{}\) ha una dipendenza implicita dall'assegnamento, e quindi sarebbe più corretto descriverla come una funzione \(\Val{}: \Ass{\phi} \to (\PL \to \set{0,1})\). Nel pratico però, l'assegnamento sarà fissato a priori.}:
  \[\Val{\alpha}[\phi] = \begin{cases}
  \alpha(\phi) &, \phi \in \AP[\phi] \\
  1 - \Val{\alpha}(\phi') &, \phi = \neg \phi' \\
  \min\set{\Val{\alpha}[\phi_1], \Val{\alpha}[\phi_2]} &, \phi = \phi_1 \land \phi_2 \\
  \max\set{\Val{\alpha}[\phi_1], \Val{\alpha}[\phi_2]} &, \phi = \phi_1 \lor \phi_2 \\
  1 &, \phi = \phi_1 \rightarrow \phi_2 \text{ e } \Val{\alpha}[\phi_1] \leq \Val{\alpha}[\phi_2] \\
  0 &, \phi = \phi_1 \rightarrow \phi_2 \text{ e } \Val{\alpha}[\phi_1] > \Val{\alpha}[\phi_2]
  \end{cases}\]
\end{definition}

A questo punto, verificare se una certa formula assume valore \(0\) o \(1\) rispetto a un assegnamento compatibile diventa un problema facilmente decidibile. 
Infatti, la valutazione di una formula \(\phi\) ha complessità \(\BigTheta{\abs{\phi}}\).

Per poter dire di aver risolto definitivamente il problema dell'infinità dei modelli, bisogna dimostrare che questa nuova semantica, basata sugli assegnamenti, sia equivalente a quella basata su modelli. 
In altre parole, dobbiamo mostrare che a ogni modello corrisponde un assegnamento equivalente, e viceversa, e che la relazione di soddisfacimento \(m \models \phi\) è equivalente alla valutazione \(\Val{\alpha}[\phi] = 1\).
Per fare ciò, andremo a mostrare che una formula è tautologica se e solo se è valida. 
Facendo ciò, otterremo anche un modo per correlare \(m \models \phi\) e \(\Val{\alpha}[\phi] = 1\).
Prima, però, definiamo il concetto di \keyword{tautologia}.

\begin{definition}{Tautologia}
  Una formula \(\phi\) è una \keyword{tautologia} (\(\Taut[\phi]\)) quando, per ogni assegnamento compatibile \(\alpha\), si ha \(\Val{\alpha}[\phi]=1\).
\end{definition}

\begin{theorem}[label=thm:validity-equals-taut]{Equivalenza tra validità e tautologicità}
  Data \(\phi \in \PL\), \(\models \phi \) se e solo se \(\Taut[\phi]\).
\end{theorem}

In precedenza, abbiamo osservato che \(\models \phi \iff \forall m, m \models \phi\).
Per definizione di tautologia, si ha che
\(\Taut[\phi]\iff\forall \alpha ,\Val{\alpha}[\phi] = 1\).
In altre parole, per dimostrare il teorema, basta mostrare che \(\forall m, m \models \phi \iff \forall \alpha, \Val{\alpha}[\phi] =1\).

\begin{lemma}[listtitle={Estrapolato dalla dimostrazione del \Cref{thm:validity-equals-taut}}]{\infosymbol{} Estrapolato dalla dimostrazione del \Crefalt[yellow]{thm:validity-equals-taut}}
  Dato un modello \(m \in 2^{AP}\), esiste un \(\alpha \in Ass(\phi)\) tale per cui \(m \models \phi \iff \Val{\alpha}[\phi]=1\).
\end{lemma}

\begin{proof}
  Dato \(m\), possiamo definire un assegnamento \(\alpha_{m}\) tale che
  \[\alpha_{m}(P) = \begin{cases}
    1, & P \in m\\
    0, & \text{ altrimenti}
  \end{cases}\]
  Possiamo mostrare quindi, per induzione strutturale, che \(m\models \phi \iff \Val{\alpha_m}[\phi] = 1\).
  \begin{description}
    \item[\(\phi \in \AP\):]
      \[m \models \phi \iff \phi \in m \eqdef {\iff}\alpha_m(\phi) = 1 \iff \Val{\alpha_m}[\phi] = 1\]
    \item[\(\phi = \neg \phi'\):]
      \[m \models \phi \iff m \not\models \phi' \stackrel{ind. hp.}{\iff} \Val{\alpha_m}[\phi'] = 0\]
      Ma allora dalla definizione di valutazione, \(\Val{\alpha_m}[\phi] = 1-\Val{\alpha_m}[\phi'] = 1-0 = 1\).
    \item[\(\phi = \phi_1 \land \phi_2\):]
      \[m \models \phi \iff m \models \phi_1 \text{ e } m \models \phi_2 \overset{ind. hp.}{\iff} \Val{\alpha_m}[\phi_1] = 1 \text{ e } \Val{\alpha_m}[\phi_2] = 1\]
      Ma allora dalla definizione di valutazione, \(\Val{\alpha_m}[\phi] = \min\set{\Val{\alpha_m}[\phi_1], \Val{\alpha_m}[\phi_2]} = \min\set{1, 1} = 1\).
    \item[\(\phi = \phi_1 \lor \phi_2\):]
      \[m \models \phi \iff m \models \phi_1 \text{ o } m \models \phi_2 \overset{ind. hp.}{\iff} \Val{\alpha_m}[\phi_1] = 1 \text{ o } \Val{\alpha_m}[\phi_2] = 1\]
      Ma allora dalla definizione di valutazione, \(\Val{\alpha_m}[\phi] = \max\set{\Val{\alpha_m}[\phi_1], \Val{\alpha_m}[\phi_2]} = \max\set{1, 0} = 1\).
    \item[\(\phi = \phi_1 \rightarrow \phi_2\):] 
      \[m \models \phi \iff \text{ se } m \models \phi_1 \text{ allora } m \models \phi_2\]
      Per ipotesi induttiva, se \(m \not\models \phi_1\) allora \(\Val{\alpha_m}[\phi_1] = 0\), e quindi a prescindere dal valore di \(\Val{\alpha_m}[\phi_2]\) si ha \(\Val{\alpha_m}[\phi] = 1\), essendo che \(0 \leq \Val{\alpha_m}[\phi_2]\).
      Se invece \(m \models \phi_1\) allora \(\Val{\alpha_m}[\phi_1] = 1\), ma sempre per ipotesi induttiva \(m \models \phi_2\) e quindi \(\Val{\alpha_m}[\phi_2] = 1\), e dunque \(\Val{\alpha_m}[\phi] = 1\), essendo che \(1 \leq 1\).
  \end{description}
\end{proof}

\begin{note}{}
  Tale dimostrazione e quella della \Namedcref{prop:subformula-upperbound} sebbene entrambe basate sull'induzione strutturale, sono molto diverse tra loro. Nella dimostrazione del \Namedonlyhref{prop:subformula-upperbound} infatti è stato possibile \qi{accorpare} i casi di operatori binari, mentre nella dimostrazione del lemma è stato necessario trattare separatamente ogni connettivo, a causa della diversa semantica di ciascuno di essi.

  Questo perché la proposizione riguardava una proprietà \textbf{sintattica} delle formule, per la quale ogni simbolo non ha altra caratterizzazione che la sua presenza nella grammatica, mentre il lemma riguarda una proprietà \textbf{semantica} delle formule, per la quale ogni simbolo ha una caratterizzazione specifica che deve essere rispettata.

  Tale differenza tra proprietà sintattiche e semantiche è fondamentale per comprendere la natura dei risultati che andremo a dimostrare, oltre a essere un fondamentale aiuto per prevedere che tipo di dimostrazione sarà necessario per dimostrare un certo risultato.
\end{note}

\begin{lemma}[listtitle={Estrapolato dalla dimostrazione del \Cref{thm:validity-equals-taut}}]{\infosymbol{} Estrapolato dalla dimostrazione del \Crefalt[yellow]{thm:validity-equals-taut}}
  Dato un assegnamento \(\alpha \in Ass(\phi)\), esiste un \(m \in 2^{AP}\) tale per cui \(m \models \phi \iff \Val{\alpha}[\phi]=1\).
\end{lemma}

\begin{proof}
  Dato un assegnamento \(\alpha\), è possibile definire un modello \(m_{\alpha}\) t.c.
  \[m_\alpha = \set{P \in \AP[\phi]}[\alpha(P) = 1]\]
  Si dimostra per induzione strutturale, in maniera del tutto analoga alla dimostrazione del precedente lemma, che \(\Val{\alpha}[\phi] = 1 \iff m_\alpha \models \phi\).
\end{proof}

Procediamo adesso, per contrapposizione, alla dimostrazione del teorema, ovvero a dimostrare \(\not\models \phi \iff \neg \Taut(\phi)\).
\begin{proof}[Dimostrazione (\ref{thm:validity-equals-taut})]
  \begin{description}
    \item[\qi{\(\implies\)}]
      Sia \(m\) tale che \(m \models \neg \phi\).
      Costruendo \(\alpha_{m}\) secondo il lemma, si ha:
      \[m\models \neg \phi \iff \Val{\alpha_{m}}[\neg\phi]=1 \iff \Val{\alpha_{m}}[\phi]=0\]
      Allora, abbiamo trovato un assegnamento \(\alpha_{m}\) tale per cui la valutazione di \(\phi\) è \(0\). Quindi, \(\neg\Taut[\phi]\).
    \item[\qi{\(\impliedby\)}]
      Sia \(\alpha\) un assegnamento per \(\phi\) tale che \(\Val{\alpha}[\phi]=0\).
      Costruiamo \(m_{\alpha}\) secondo il lemma.
      Allora, si ha:
      \[\Val{\alpha}[\neg\phi]=1 \iff m \not\models \phi\]
      Allora, abbiamo trovato un modello \(m_{\alpha}\) che non soddisfa \(\phi\). Quindi, \(\not\models \phi\).
  \end{description}
\end{proof}

Questo teorema ci dice che risolvere \(\Taut\) equivale a risolvere \(\Valid\) e che, quindi, essendo \(\Taut\) computabile, in quanto una ricerca in un insieme finito, anche \(\Valid\) lo è.
Il costo computazionale di \(\Valid\) è però esponenziale, in quanto \(\abs{\Ass{\phi}} = 2^{\abs{\phi}}\) e dunque \(\Valid\) ha complessità \(\BigO{\abs{\phi}2^{\abs{\phi}}}\leq 2^{\BigTheta{\abs{\phi}}}\).

\begin{example}[namedlabel={ex:notable_tautologies}{Tautologie notevoli}]{Tautologie notevoli}
  Vediamo alcune tautologie notevoli che ci saranno successivamente utili per la dimostrazione di altri risultati.
  \begin{enumerate}
    \item \(P \lor \neg P\) (legge del terzo escluso).
    \item \(\neg (P \land \neg P)\) (principio di non contraddizione \(\bot\))
    \item \(P \lor Q \leftrightarrow \neg(\neg P \land \neg Q)\) (legge di De Morgan)
    \item \(P\land Q \leftrightarrow \neg(\neg P \lor \neg Q)\) (legge di De Morgan)
    \item \(\neg \neg P \iff P\) (eliminazione della doppia negazione)
    \item \(P \land(Q \lor R) \leftrightarrow (P \land Q) \lor (P \land R)\) (distributività)
    \item \((P \lor(Q\land R)) \leftrightarrow (P \lor Q) \land (P \lor R)\) (distributività)
    \item \((P \rightarrow Q) \leftrightarrow (\neg P \lor Q)\) (implicazione materiale)
  \end{enumerate}
\end{example}

\section{Sostituzioni nella Logica Proposizionale}

Per avvicinarci gradualmente verso la formalizzazione del concetto di \q{ragionamento}, o deduzione, anticipato all'inizio del corso, introduciamo una prima forma di manipolazione delle formule: la sostituzione, che permette di costruire nuove formule a partire da formule già esistenti, preservando alcune proprietà semantiche.

\begin{definition}{Sostituzione}
  Definiamo una \keyword{sostituzione} come una funzione \(\sigma: \AP \to \PL\).
\end{definition}

La sostituzione permette di sostituire ogni proposizione atomica con una formula arbitraria, e, dunque, di costruire nuove formule a partire da formule già esistenti. Procediamo, adesso, con la definizione formale.

\begin{definition}{Applicazione di una sostituzione a una formula}
  Dato \(\phi \in \PL\) e una sostituzione \(\sigma\), definiamo l'applicazione di \(\sigma\) a \(\phi\), denotata con \(\phi\sigma\), come segue:
  \begin{itemize}
   \item se \(\phi \in \AP\) allora \(\phi\sigma = \sigma(\phi)\)
   \item se \(\phi = \neg \phi'\) allora \(\phi\sigma = (\neg \phi ')\sigma = \neg (\phi'\sigma)\)
   \item se \(\phi = \phi_1 \odot \phi_2\) allora \(\phi\sigma = (\phi_1 \odot \phi_2)\sigma = (\phi_1\sigma) \odot (\phi_2\sigma)\)
  \end{itemize}
\end{definition}

\begin{note}{}
  È importante sottolineare che la sostituzione è da intendersi come applicata simultaneamente a tutte le proposizioni atomiche. Se una proposizione atomica \(P\) compare sia come simbolo da sostituire che come simbolo di una formula sostituita, la formula sostituita a \(P\) non viene ulteriormente sostituita anche internamente alle formule sostituite contenenti \(P\).
\end{note}

\begin{example}{Sostituzioni e loro applicazione}  
  Consideriamo \(\sigma = \set{A \leftarrow (A\lor B), B \leftarrow \neg A, C \leftarrow C}\) e \(\phi = (A \land B) \lor C\).
  
  Allora, \(\phi\sigma\) sarà:
  \[((A\lor B) \land \neg A) \lor C\]

  e non:
  \[((A\lor \neg A)\land \neg A)\lor C\]

  Nel secondo caso infatti si è applicata la sostituzione per \(B\) anche all'interno della formula sostituita a \(A\). Tale comportamento è quello che si sarebbe ottenuto considerando un ulteriore sostituzione \(\tau = \set{B\leftarrow \neg A}\) e applicando \((\phi\sigma)\tau\).
\end{example}

Per mostrare come la sostituzione sia una prima forma di deduzione logica molto elementare, e che quindi preservi alcune proprietà semantiche, andiamo a enunciare e successivamente dimostrare il \Namedonly{thm:substitution}.

\begin{theorem}[namedlabel={thm:substitution}{Teorema di Sostituzione}]{Teorema di Sostituzione}
  Sia \(\phi\in \Taut\) e \(\sigma\) un'arbitraria sostituzione. Allora \(\phi\sigma\in \Taut\).
\end{theorem}

Per dimostrare questo teorema, è necessario dare carattere semantico all'operazione di sostituzione, altrimenti esclusivamente sintattica, e mostrare che questa operazione preserva la proprietà di essere una tautologia.

\begin{definition}{Semantica della sostituzione}
  Data \(\phi \in \PL\), una sostituzione \(\sigma\) e un assegnamento \(\alpha \in \Ass{\phi\sigma}\), definiamo l'assegnamento \(\alpha^\sigma \in \Ass{\phi}\) come:
  \[\alpha^\sigma(P)=\Val{\alpha}[\sigma(P)]\]
\end{definition}

In altre parole, \(\alpha^{\sigma}\) è un assegnamento che, a ciascuna proposizione elementare \(P\in Sub(\phi)\), assegna la valutazione della formula \(\sigma(P)\) ottenuta dopo la sostituzione \(\sigma\).

\begin{note}{}
  Si noti che \(\alpha^\sigma\) è un assegnamento per \(\phi\), mentre \(\alpha\) è un assegnamento per \(\phi\sigma\).
\end{note}

Vediamo dunque un risultato intermedio che ci sarà utile per dimostrare il \Namedonlyhref{thm:substitution}.
\begin{lemma}[namedlabel={lem:substitution}{Lemma di Sostituzione}]{Lemma di Sostituzione}
  \(\forall \phi \in \PL, \sigma \in \PL^{\AP}, \alpha \in \Ass{\phi\sigma}\) vale che
  \[\Val{\alpha}[\phi\sigma]=\Val{\alpha^\sigma}[\phi]\]
\end{lemma}

\begin{proof}
  Dimostriamo per induzione strutturale.

  \begin{description}
    \item[\(\phi \in \AP\):] in questo caso, che rappresenta il caso base, \(\phi\sigma = \sigma(\phi)\). Ma quindi
      \[\Val{\alpha}[\phi\sigma] = \Val{\alpha}[\sigma(\phi)] \eqdef \alpha^\sigma(\phi) = \Val{\alpha^\sigma}[\phi]\]
    \item[\(\phi = \neg \psi\):] in questo primo caso induttivo, dalla definizione di sostituzione abbiamo che \(\phi\sigma = \neg (\psi\sigma)\). Allora:
      \[\Val{\alpha}[\phi\sigma] = \Val{\alpha}[\neg (\psi\sigma)] = 1 - \Val{\alpha}[\psi\sigma]\]
      Ma per ipotesi induttiva \(\Val{\alpha}[\psi\sigma] = \Val{\alpha^\sigma}[\psi]\), quindi
      \[\Val{\alpha}[\phi\sigma] = 1 - \Val{\alpha^\sigma}[\psi] = \Val{\alpha^\sigma}[\neg \psi] = \Val{\alpha^\sigma}[\phi]\]
    \item[\(\phi = \psi_1 \odot \psi_2\):] Per tutti gli operatori binari, che rappresentano i restanti casi induttivi, la struttura di dimostrazione è analoga, e si basa sempre sull'ipotesi induttiva. Ad esempio, se \(\phi = \psi_1 \land \psi_2\) allora
      \[\Val{\alpha}[\phi\sigma] = \Val{\alpha}[(\psi_1\sigma) \land (\psi_2\sigma)] = \min\set{\Val{\alpha}[\psi_1\sigma], \Val{\alpha}[\psi_2\sigma]}\]
      Ma per ipotesi induttiva \(\Val{\alpha}[\psi_i\sigma] = \Val{\alpha^\sigma}[\psi_i]\) per \(i\in\set{1,2}\), quindi
      \[\Val{\alpha}[\phi\sigma] = \min\set{\Val{\alpha^\sigma}[\psi_1], \Val{\alpha^\sigma}[\psi_2]} = \Val{\alpha^\sigma}[\psi_1 \land \psi_2] = \Val{\alpha^\sigma}[\phi]\]
  \end{description}
\end{proof}

\begin{note}{}
  Si noti che questa dimostrazione non contraddice la nota precedente riguardo alla differenza tra proprietà sintattiche e semantiche, in quanto nonostante per brevità siano state omesse le dimostrazioni per i diversi operatori binari, esse sono comunque diverse tra loro, in quanto ogni connettivo ha una semantica diversa che deve essere rispettata. Per mostrare infatti la dimostrazione di tale caso è stato necessario istanziare \(\odot\) con un operatore specifico, e mostrare che la semantica di tale operatore è preservata dalla sostituzione.
\end{note}

Forti del lemma, possiamo finalmente procedere con la dimostrazione del \Namedonlyhref{thm:substitution}, che ora risulta quasi immediata.
Vogliamo dimostrare che \(\forall \sigma \in \PL^{\AP}\), \(\phi \in \Taut \implies \phi\sigma \in \Taut\).
Procederemo per contrapposizione, ovvero assumeremo che \(\phi\sigma \not\in \Taut\) e dimostreremo che \(\phi \not\in \Taut\).

\begin{proof}[Dimostrazione del (\ref{thm:substitution})]
  Se \(\phi\sigma \not\in \Taut\), allora \(\exists \alpha \in \Ass{\phi\sigma} \st \Val{\alpha}[\phi\sigma] = 0\).

  Per il \Namedonlyhref{lem:substitution}, \(\Val{\alpha}[\phi\sigma] = \Val{\alpha^\sigma}[\phi]\), quindi \(\Val{\alpha^\sigma}[\phi] = 0\), ovvero esiste un assegnamento \(\alpha^\sigma\) tale che \(\phi\) è falsa, e quindi \(\phi \not\in \Taut\).
\end{proof}


\begin{observation}{}
  È importante notare che, nonostante la sostituzione preservi validità e tautologicità, essa non preserva la soddisfacibilità o, in generale, la verità.
  Formalmente:
  \[\phi \in \Sat \centernot\implies \phi\sigma \in \Sat\]

  Questo è facile da vedere considerando \(\phi = P \in \AP\). Si osserva \(\phi \in \Sat\) per qualsiasi assegnamento \(\alpha\) tale per cui \(\alpha(P)=1\).
  Scegliendo però \(\sigma : P \to P\land \neg P\), che è chiaramente insoddisfacibile, e fissato un modello \(m\), si ha che
  \[m\models\phi \centernot\implies m\models\phi\sigma\]
  Altrimenti, si otterrebbe una contraddizione.

  Invece, la sostituzione conserva l'insoddisfacibilità. Infatti,
  \[\phi \in \Unsat \iff \neg \phi \in \Taut \implies (\neg\phi)\sigma \in \Taut\iff \]
  \[\iff \neg(\phi\sigma) \in \Taut \implies \phi\sigma \in \Unsat\]
\end{observation}

Il teorema di sostituzione è un primo esempio di come sia possibile manipolare le formule in modo da preservare alcune proprietà semantiche. 
In particolare, tale teorema ci dice che è possibile creare infinite formule valide, poiché è sufficiente partire da una formula valida, applicare una qualsiasi sostituzione e ottenere una nuova formula valida.

\subsection{Deduzione Semantica}

Possiamo notare che \(A \to (B\to C)\equiv (A\land B)\to C\).
Infatti, chiamando \(\phi_1 = A \to (B\to C)\) e la seconda \(\phi_2 = (A\land B)\to C\) possiamo andare ad analizzare quali assegnamenti rendono vere o false ciascuna di queste formule.

La prima formula \(\phi_1\) ha come connettivo principale l'implicazione, e dunque è falsa se e soltanto se \(A\) è vera e \(B\to C\) è falsa. Ma \(B\to C\) è un altra implicazione, dunque è falsa se e soltanto se \(B\) è vera e \(C\) è falsa. Di conseguenza, \(\phi_1\) è falsa se e soltanto se \(A\) è vera, \(B\) è vera e \(C\) è falsa.

D'altro canto, anche la seconda formula \(\phi_2\) ha come connettivo principale l'implicazione, e dunque è falsa se e soltanto se \(A\land B\) è vera e \(C\) è falsa. Ma \(A\land B\) è vera se e soltanto se \(A\) è vera e \(B\) è vera. Di conseguenza, \(\phi_2\) è falsa se e soltanto se \(A\) è vera, \(B\) è vera e \(C\) è falsa, ovvero le due formule sono false, e quindi vere, negli stessi assegnamenti, e dunque sono logicamente equivalenti.

Possiamo dunque trasformare questa equivalenza in una tautologia, scrivendo
\[A \to (B\to C) \leftrightarrow (A\land B)\to C\]

Questa tautologia è alla base del concetto di deduzione semantica, ovvero la possibilità di dedurre una formula \(\phi\) sapendo che esiste una formula \(\psi\) vera tale che \(\psi \to \phi\) è vera.

\begin{theorem}[namedlabel={thm:semantic_deduction}{Teorema di Deduzione Semantica}]{Teorema di Deduzione Semantica}
  Sia \(\Gamma \subseteq \PL\) e \(\phi,\psi \in \PL\). Allora
  \[\Gamma,\phi \models \psi \iff \Gamma \models \phi\to\psi\]
\end{theorem}

Il teorema permette di \qi{estrarre} una premessa, spostandola nella conseguenza logica da dimostrare.
Ciò consente spesso di semplificare le dimostrazioni. 
Infatti, per dimostrare \(\Gamma,\phi \models \psi\), è necessario, in primo luogo, trovare un modello per cui valga \(\Gamma, \phi\).
Abbiamo visto che la ricerca di un modello ha complessità esponenziale rispetto al numero di proposizioni elementari coinvolte. A ciò, si aggiungono eventuali dipendenze tra queste. 
Pertanto, cercare un modello che soddisfi \(\Gamma\) è sicuramente più facile che cercare un modello che soddisfi \(\Gamma, \phi\).
Inoltre, la deduzione logica permette di soddisfare le proposizioni in maniera sequenziale.
Ad esempio, invece di dire \qi{Se piove e fa freddo, allora prendo l'ombrello e la sciarpa}, possiamo \qi{Se piove, allora se fa freddo prendo l'ombrello e la sciarpa}, il che ci consente di affrontare un problema alla volta.

\begin{proof}
  Dalla definizione di conseguenza logica, \(\Gamma,\phi \models \psi\) significa che
  \[\forall m \in 2^{\AP}, m\models \Gamma \cup \set{\phi} \implies m\models \psi\]

  L'argomento sinistro dell'implicazione è però equivalente a dire che \(m\models \Gamma\) e \(m\models \phi\). Quindi
  \[\forall m \in 2^{\AP}, m\models \Gamma \text{ e } m\models \phi \implies m\models \psi\]

  Dando quindi semantica agli operatori di metalivello nel modo usuale (ovvero identico a quelli della logica proposizionale), possiamo riscrivere la precedente formula con variabili proposizionali \(A\), \(B\) e \(C\) al posto di \(m\models \Gamma\), \(m\models \phi\) e \(m\models \psi\) rispettivamente, ottenendo
  \[\forall m \in 2^{\AP}, A\text{ e }B \implies C\]
  che, per la tautologia vista in precedenza, è equivalente a
  \[\forall m \in 2^{\AP}, A \implies (B \implies C)\]
  Applicando la sostituzione \(\sigma=\set{ A \to m\models \Gamma, B \to m\models \phi, C \to m\models \psi }\), si ottiene 
  \[\forall m \in 2^{\AP}, m\models \Gamma \implies (m\models \phi \implies m\models \psi)\]

  Per semantica dell'implicazione si ha che \((m\models \phi \implies m\models \psi) \iff (m\models \phi \to \psi)\), dunque
  \[\forall m \in 2^{\AP}, m\models \Gamma \implies m\models \phi\to\psi\]
\end{proof}

\section{Forme normali}

Il linguaggio che abbiamo utilizzato fino ad adesso contiene molti simboli, ma come abbiamo già osservato tramite alcune tautologie è possibile esprimere alcuni operatori in funzione di altri.

Un esempio potrebbero essere le leggi di De Morgan
\[\neg(\phi\land\psi) \equiv \neg\phi \lor \neg\psi\]
\[\neg(\phi\lor\psi) \equiv \neg\phi \land \neg\psi\]

Interpretando queste equivalenze come regole di trasformazione, questo ci permette di spingere le negazioni fino alle proposizioni atomiche mantenendo l'equivalenza logica.

In questo modo, possiamo definire una prima forma normale, chiamata \(\PNF\) (Positive Normal Form) o \(\NNF\) (Negation Normal Form), che contiene tutte le formule di \(\PL\) in cui le negazioni occorrono solo davanti a proposizioni atomiche.

Per definire formalmente questa forma normale, introduciamo il concetto di letterale.

\begin{definition}{Letterale}
  Una formula \(\phi\) è un \keyword{letterale} se \(\phi \in AP\) o \(\phi = \neg P \), con \(P \in AP\).
\end{definition}

\begin{definition}{Forma normale positiva}
  Definiamo la forma normale positiva (\(\PNF\) o \(\NNF\)) come l'insieme di tutte le formule costruite come in \(\PL\), senza usare negazioni e utilizzando l'insieme dei letterali \(\mathcal{L}\) come \(\AP\).
\end{definition}

Si ha ovviamente che \(\PNF \subset \PL\), ma dalle leggi di De Morgan è ovvio che
\[\forall \phi \in \PL, \exists \phi' \in \PNF: \phi \iff \phi'\]

In generale, il concetto di forma normale è un sottoinsieme di formule, costruibile in maniera induttiva, che mantiene la completezza semantica.
Ciò significa che per ogni formula esiste una formula equivalente in forma normale.

Le due forme normali più importanti sono la \keyword{forma normale congiuntiva} (\(\CNF\)) e la \keyword{forma normale disgiuntiva} (\(\DNF\)).

\begin{definition}{Disjunctive Normal Form}
  Definiamo la \keyword{forma normale disgiuntiva} (\(\DNF\)) come l'insieme di tutte le formule con la seguente struttura:
  \[\phi = \bigvee_{i=1}^k\left(\bigwedge_{j=1}^{r_i} l_{ij}\right)\]
  dove \(l_{ij} \in \mathcal{L}\) per ogni \(i\) e \(j\).
\end{definition}

\begin{example}{}
  Un esempio di formula in DNF è la seguente:
  \[(P\land \neg Q) \lor (\neg P \land \neg Q \land R) \lor \ldots\]

  Ognuna delle parentesi (quindi ogni congiunzione) è detta \keyword{clausola congiuntiva}.
\end{example}

\begin{definition}{Conjunctive Normal Form}
  Definiamo la \keyword{forma normale congiuntiva} (\(\CNF\)) come l'insieme di tutte le formule con la seguente struttura:
  \[\phi = \bigwedge_{i=1}^k\left(\bigvee_{j=1}^{r_i} l_{ij}\right)\]
  dove \(l_{ij} \in \mathcal{L}\) per ogni \(i\) e \(j\).
\end{definition}

\begin{example}{}
  Un esempio di formula in CNF è la seguente:
  \[(P\lor \neg Q) \land (\neg P \lor \neg Q \lor R) \land \ldots\]

  Ognuna delle parentesi (quindi ogni disgiunzione) è detta \keyword{clausola disgiuntiva}.
\end{example}

Essendoci dualità tra congiunzione e disgiunzione, è facile vedere che le due forme normali sono duali.

Ovviamente \(\CNF \cup \DNF \subseteq \PNF\). È possibile mostrare facilmente come, per ogni \(\phi \in \Sat\), sia possibile costruire \(\psi \in \DNF\) tale che \(\phi \equiv \psi\) partendo dalla tabella di verità di \(\phi\).
Infatti, basterà aggiungere una clausola congiuntiva per ogni assegnamento che rende vera \(\phi\).
In ciascuna clausola congiuntiva, un \(P\in \AP[\phi]\) corrisponderà al letterale \(P\) se \(P\) è vero nell'assegnamento, al letterale \(\neg P\) altrimenti.
Il seguente esempio chiarirà sicuramente il procedimento descritto.

\begin{example}{}
  Consideriamo la formula \(\phi \) con la seguente tabella di verità:

  \begin{center}
    \begin{tabular}{ccc|c}
      \(A\) & \(B\) & \(C\) & \(\phi\)\\
      \hline
      0 & 0 & 0 & 1\\
      0 & 0 & 1 & 0\\
      0 & 1 & 0 & 1\\
      0 & 1 & 1 & 0\\
      1 & 0 & 0 & 1\\
      1 & 0 & 1 & 0\\
      1 & 1 & 0 & 1\\
      1 & 1 & 1 & 0\\
    \end{tabular}
  \end{center}

  Dalla tabella di verità, possiamo costruire la formula in DNF equivalente a \(\phi\) come segue:
  \[(\neg A \land \neg B \land \neg C) \lor (\neg A \land B \land \neg C) \lor (A \land \neg B \land \neg C) \lor (A \land B \land \neg C)\]
  
\end{example}

\begin{note}{}
  Il procedimento mostrato non permette di costruire formule non soddisfacibili. Se non riuscissimo a esprimere le formule insoddisfacibili in \(\DNF\), allora la \(\DNF\) non manterrebbe, come abbiamo detto, la completezza semantica.
  Per risolvere il \qi{buco} lasciato dal procedimento descritto, è sufficiente associare, a ogni formula insoddisfacibile, la formula \(P\land \neg P\), che è \(\DNF\).
  Ciò è possibile dal momento in cui tutte le formule insoddisfacibili sono equivalenti a una contraddizione.
\end{note}

Per le \(\CNF\) il ragionamento è simile. 
In questo caso, però, poiché dobbiamo esprimere la formula come una congiunzione di clausole disgiuntive, non possiamo semplicemente elencare degli assegnamenti come in precedenza, poiché gli assegnamenti sono essenzialmente delle congiunzioni.
Dunque, al posto di elencare gli assegnamenti che rendono vera la formula, elenchiamo in congiunzione la negazione di ogni assegnamento che la rende falsa. 
In questo modo, ci troviamo con una formula della forma:
\[\phi' = \neg(l_1^1 \land \ldots \land \ldots \land l_1^{r_1})\land\neg(l_2^{1}\land \ldots \land \ldots\land l_2^{r_2})\land\ldots\land \neg(l_k^1\land\ldots\land l_k^{r_k})\]
A questo punto, applicando le leggi di De Morgan, possiamo ottenere una formula della forma:
\[\phi'' = (\neg l_1^1 \lor \ldots \lor \neg l_1^{r_1})\land(\neg l_2^1 \lor \ldots \lor \neg l_2^{r_2})\land\ldots\land (\neg l_k^1 \lor \ldots \lor \neg l_k^{r_k})\]
che è esattamente una formula in \(\CNF\).
Si noti che, se, prima dell'applicazione della legge di De Morgan, un letterale era della forma \(\neg l\), allora l'applicazione di De Morgan porta il letterale a essere \(\neg \neg l\), che, per doppia negazione, equivale a \(l\).

\begin{example}{}
  Riprendendo la tabella di verità dell'esempio precedente, possiamo costruire la formula in \(\CNF\) equivalente a \(\phi\) come segue:
  \[\neg(\neg A \land \neg B \land C) \land \neg(\neg A \land B \land C) \land \neg(A \land \neg B \land C) \land \neg(A \land B \land C)\]
  \[\equiv (A \lor B \lor \neg C) \land (A \lor \neg B \lor \neg C) \land (\neg A \lor B \lor \neg C) \land (\neg A \lor \neg B \lor \neg C)\]
\end{example}

\begin{note}{}
  In modo duale alle \(\DNF\), il procedimento descritto non consente di costruire una formula \(\CNF\) equivalente a formule valide.
  Poiché tutte le formule valide sono equivalenti a una tautologia, è possibile associare a esse la formula \(P \lor \neg P\), che è \(\CNF\).
\end{note}

Tuttavia, queste costruzioni sono estremamente inefficienti e portano a formule di dimensione esponenziale rispetto alla formula di partenza. 
Ciò avviene perché la dimensione della tabella di verità di una formula \(\phi\) è \(2^{\abs{\AP[\phi]}}\), dovendo considerare tutti gli assegnamenti possibili, e dunque la formula in \(\DNF\) o \(\CNF\) costruita a partire da tale tabella può avere un numero di clausole esponenziale. 
In particolare, se \(\phi\) è già in \(\DNF\) o in \(\CNF\), allora passare alla forma normale duale usando le leggi di De Morgan e la distributività porta a una formula di dimensione esponenziale.

\begin{example}{}
  Sia \(\phi \in\DNF\) tale che \(\phi\) sia:
  \[(P_1 \land P_2)\lor (P_3 \land P_4)\]

  Andiamo quindi a trasformarla in \(\CNF\) usando la proprietà distributiva.
  \[(P_1 \land P_2)\lor (P_3 \land P_4)\equiv\]
  \[\equiv ((P_1 \land P_2)\lor P_3)\land((P_1 \lor P_4)\lor P_4)\equiv\]
  \[\equiv (P_1 \lor P_3) \land (P_2 \lor P_3) \land (P_1 \lor P_4) \land (P_2 \lor P_4)\]

  In questo caso si è passati da una formula con \(2\) clausole a una formula con \(4\) clausole, che non è un evidenza sufficiente per la crescita esponenziale.

  Consideriamo però una formula:
  \[\phi = (l_1^1 \land l_2^1) \lor (l_1^2 \land l_2^2) \lor\ldots\lor (l_1^n \land l_2^n)\]
  che è in \(\DNF\) e contiene \(n\) clausole congiuntive binarie.

  Analizzando la distributività però, possiamo notare come venga generata una clausola disgiuntiva per ogni possibile combinazione di letterali.

  In generale, quindi, ogni clausola disgiuntiva prodotta con la distributività corrisponde a una scelta di un letterale per ogni clausola congiuntiva, ed è dunque equivalente a una funzione \(f:\set{1,2,\ldots,n} \to \{1,2\}\). Dunque, il numero di clausole disgiuntive generate è pari al numero di tali funzioni, che è \(\abs{\set{1,2}}^n = 2^n\).
\end{example}

Il motivo per cui abbiamo introdotto queste forme è che, se una formula è in \(\DNF\), la sua soddisfacibilità è facilmente verificabile. 
Dualmente, se una formula è in \(\CNF\), la sua validità è facilmente verificabile.
Infatti, sia
\[\phi = C_1 \lor C_2 \lor \ldots \lor C_k\]
con \(C_i = (l_1^i \land \ldots \land l_{r_i}^i)\), per ogni \(i \in \set{1,\ldots,k}\).
Allora \(\phi\) è soddisfacibile se e solo se almeno una clausola congiuntiva \(C_i\) è soddisfacibile.
Inoltre, una clausola congiuntiva \(C_i\) è insoddisfacibile se e solo se contiene due letterali complementari, cioè esistono \(l_j^i\) e \(l_h^i\) tali che \(l_j^i = \neg l_h^i\).
Infatti, in presenza di una coppia complementare compare una contraddizione del tipo \(P \land \neg P\); viceversa, se non compaiono letterali complementari, è sempre possibile costruire un assegnamento che rende vera la clausola.

Dualmente, sia
\[\phi = D_1 \land D_2 \land \ldots \land D_k\]
con \(D_i = (l_1^i \lor \ldots \lor l_{r_i}^i)\), per ogni \(i \in \set{1,\ldots,k}\).
Allora \(\phi\) è valida se e solo se tutte le clausole disgiuntive \(D_i\) sono valide.
Inoltre, una clausola disgiuntiva \(D_i\) è valida se e solo se contiene due letterali complementari, cioè esistono \(l_j^i\) e \(l_h^i\) tali che \(l_j^i = \neg l_h^i\).
